\documentclass[11pt,a4paper]{article}

\usepackage{amsmath,amssymb,amsthm}
\usepackage{lmodern}
\usepackage[T1]{fontenc}
\usepackage[utf8]{inputenc}
\usepackage{hyperref}
\usepackage{xcolor}
\usepackage{geometry}
\geometry{margin=1in}

\hypersetup{
  pdftitle={The Riemann Hypothesis via a Bounded-Coefficient Prior and Regular Conditional Probability},
  pdfauthor={Leonardo Pedro},
  colorlinks=true,
  linkcolor=blue,
  citecolor=blue,
  urlcolor=blue
}

\newtheorem{theorem}{Theorem}[section]
\newtheorem{lemma}[theorem]{Lemma}
\newtheorem{proposition}[theorem]{Proposition}
\newtheorem{corollary}[theorem]{Corollary}
\theoremstyle{definition}
\newtheorem{definition}[theorem]{Definition}
\newtheorem{remark}[theorem]{Remark}

\newcommand{\Rea}{\operatorname{Re}}
\newcommand{\Ima}{\operatorname{Im}}
\newcommand{\etaX}{\eta_X}
\newcommand{\rcp}{\operatorname{rcp}}

\title{The Riemann Hypothesis via a Bounded-Coefficient Prior\\and Regular Conditional Probability}
\author{Leonardo Pedro}
\date{}

\begin{document}
\maketitle

\begin{abstract}
We describe, and partly formalize in Lean~4, a strategy for the Riemann Hypothesis
that works \emph{directly} in standard complex analysis (ZFC), with no arithmetical
encoding. To the deterministic Dirichlet \(\eta\)-function we attach a \emph{random}
Dirichlet series \(\etaX(s)=\sum_n X_n n^{-s}\) whose coefficients are independent,
\textbf{uniformly bounded} (\(\|X_n\|\le 2\)) --- critical, as this is what gives
absolute convergence on \(\Rea(s)>1\) and anchors the random series to \(\eta\) ---
and have \textbf{mean} \(\mathbb{E}[X_n]=1\), so the mean realization \(X_n\equiv1\) is
exactly \(\eta\). The coefficients need \emph{not} be multiplicative: there is no Euler
product. The prior is the atomless uniform law on the closed disk of \textbf{radius
\(2\)}, chosen so that the mean realization \(X_n\equiv1\) lies in the \emph{interior}
of the support, where a regular conditional probability (rcp) is unambiguously well
defined.

The single role of this Mehler-type probability measure is to supply one identity --- the
\textbf{transfer equality}:
\[
  \mathbb{P}\bigl(\text{standard }\eta\text{ has a zero at }s_0\bigr)
  \;=\;
  \rcp\bigl(\etaX\text{ has a zero at }s_0 \,\big|\, X_n=1\ \forall n\bigr),
\]
the rcp at the interior mean realization of the atomless Radon measure returning the
deterministic value of the (continuous) ``zero at \(s_0\)'' functional. The left side is
a deterministic \(\{0,1\}\) indicator. The right side is then bounded \emph{away from
\(1\)} by \textbf{Bagchi/Voronin universality}: conditioning on a \emph{local} edge
neighborhood of \(\eta\) (never a contour enclosing \(s_0\), which the maximum-modulus
principle would forbid), a positive-mass family of zero-free realizations matches \(\eta\)
on the edges yet stays nonzero at \(s_0\), so \(L:=\rcp(\,\cdot\mid X_n=1)\neq 1\). A
\(\{0,1\}\) indicator equal to something \(\neq 1\) is \(0\): \(\eta\) has \emph{no} zero
at \(s_0\). Uniformly over \(1/2<\Rea(s_0)<1\) this is the Riemann Hypothesis --- a
\emph{positive} theorem in ZFC, with no \(\Pi^0_1\) encoding, no witness extraction, and
no unprovability hedge. The route has \textbf{two} deep, un-formalized analytic inputs:
the \textbf{transfer equality} (the role of the Mehler measure) and \textbf{Bagchi (1981)
universality} (non-detection).

The probabilistic substrate is an instance of the \textbf{Kopperman/Mehler formalism}
used in the companion \(P\neq NP\) development: a separable Hilbert space, a decidable
dense skeleton of finite Dirichlet approximants, and an \emph{atomless} measure specified
by its mean and covariance. ``Choosing the measure is choosing the model.'' In Lean, the
probabilistic core \texttt{RcpEuler.lean} (the random series, the atomless prior, the rcp
kernel, the reciprocal contour convergence) is on disk and \texttt{sorry}-free for an
earlier radius-\(1\), Euler-product instance; the redesign of this paper --- radius \(2\),
general non-multiplicative series, and the transfer-equality bridge in place of the
arithmetical one --- is the current plan (\texttt{IMPLEMENTATION\_PLAN\_RCP.md}) and is
not yet on disk. The result is a \emph{verified reduction} of RH to the transfer equality
and Bagchi universality, not a complete proof.
\end{abstract}

\section{Introduction}

\subsection{Strategy: make a deterministic question probabilistic}

The classical attack on RH tries to show directly that \(\zeta\) (equivalently \(\eta\))
has no zero in the open right half \(\{1/2<\Rea(s)<1\}\) of the critical strip. ``Does
\(\eta\) vanish at \(s_0\)?'' is a single deterministic question with a \(\{0,1\}\)
answer; analytically it is intractable.

The present strategy replaces it by a \emph{probabilistic} question that has the same
answer but is amenable to a soft argument. We embed \(\eta\) as the mean of a random
Dirichlet series \(\etaX(s)=\sum_n X_n n^{-s}\), \(\mathbb{E}[X_n]=1\), over an atomless
prior, and ask for the regular conditional probability that \(\etaX\) vanishes at \(s_0\)
\emph{given} that the coefficients take their mean values \(X_n=1\). Two facts then close
the argument.

\paragraph{Fact 1 --- the transfer equality (the role of the Mehler measure).} Adding the
probability measure buys exactly one identity:
\begin{equation}\label{eq:transfer}
  \mathbb{P}\bigl(\eta(s_0)=0\bigr)
  \;=\;
  \rcp\bigl(\etaX(s_0)=0 \,\big|\, X_n=1\ \forall n\bigr)
  \;=:\;L .
\end{equation}
The left side is the deterministic indicator \(\mathbf{1}\{\eta(s_0)=0\}\in\{0,1\}\): the
standard \(\eta\) either vanishes at \(s_0\) or it does not. The right side is the rcp at
the mean realization. Equality~\eqref{eq:transfer} holds because the rcp at an
\emph{interior} point of the support of an atomless Radon measure returns the
deterministic value of any functional continuous there, and ``has a zero at \(s_0\)'' is
such a functional (Section~\ref{sec:transfer}). The choice of \textbf{radius \(2\)} for
the coefficient disk is what puts \(X_n\equiv1\) in the interior, where this is clean.

\paragraph{Fact 2 --- non-detection (\(L\neq 1\), Bagchi).} Conditioning on a
\emph{local} edge neighborhood of \(\eta\) --- never a contour enclosing \(s_0\), which
the maximum-modulus principle would forbid --- \textbf{Bagchi/Voronin universality}
provides a positive-mass family of zero-free realizations that agree with \(\eta\) on the
edges but stay nonzero at \(s_0\). Hence \(L\neq 1\): the conditional law of \(\etaX(s_0)\)
is not a point mass at \(0\) (Section~\ref{sec:nondetect}).

\paragraph{Composing.} A \(\{0,1\}\)-valued indicator that equals \(L\neq 1\) must equal
\(0\). Equivalently: if \(\eta(s_0)=0\) then \eqref{eq:transfer} forces \(L=1\),
contradicting Fact~2. Either way \(\eta(s_0)\neq 0\). Since the universality argument is
uniform over the interior point,
\[
  \forall\,s_0,\ \tfrac12<\Rea(s_0)<1 \;\Longrightarrow\; \eta(s_0)\neq 0,
\]
which is the Riemann Hypothesis (Section~\ref{sec:assembly}). The conclusion is a positive
analytic statement in ZFC; there is no arithmetical sentence and no unprovability reading.

\subsection{Why ``regular conditional probability'' and not a single condition}

The conditioning event ``\(\etaX\) equals \(\eta\) on the (local) edge data'' is null: the
boundary trace pins a continuum of constraints, and the mean realization \(X_n\equiv1\) is
a single point of an atomless measure. Regular conditional probability is precisely the
tool that assigns a well-defined conditional law to such null conditioning data, via
Tjur's neighborhood-limit characterization
\begin{equation}\label{eq:tjur}
  \rcp(A \mid T = t)
  \;=\; \lim_{V \downarrow \{T=t\}} \frac{\mathbb{P}(A \cap V)}{\mathbb{P}(V)},
\end{equation}
the limit over shrinking open neighborhoods \(V\) of the trace value. The limit is well
defined because the substrate is \textbf{Radon} (a Polish inner-product space has
inner-regular Borel measures), and --- crucially --- because the conditioning value
\(t=e_\eta\) lies in the \emph{interior} of the support (the radius-\(2\) choice), so the
neighborhoods \(V\) are genuine two-sided neighborhoods with positive mass in all
directions. It is this limit, not a naive ratio, that both the transfer
equality~\eqref{eq:transfer} and the non-detection bound \(L\neq1\) control.

\section{The Bounded-Coefficient Prior and the rcp Kernel}

\subsection{The model: a general, non-multiplicative random Dirichlet series}\label{sec:model}

The object is a random Dirichlet series. There is \textbf{no Euler product and no
multiplicativity}; only boundedness and a prescribed mean and covariance are used.

\begin{definition}[Random Dirichlet series]\label{def:series}
Fix a sample \(\omega\in\Omega_b=\mathbb{N}\to\mathbb{C}\) recording one complex
coefficient per index. With \(X_n(\omega)=\omega_n\) set
\[
  \etaX(s,\omega) \;=\; \sum_{n\ge 1} X_n(\omega)\,n^{-s}
  \qquad(\eta\text{-normalized, alternating}),
\]
the random series whose mean realization \(X_n\equiv1\) is the deterministic \(\eta\). The
exponential Euler product \(\mathcal{B}_P(s,\omega)=\exp\bigl(\sum_{p\le P}X_p\,p^{-s}
+\log\eta^E_P(s)\bigr)\) of earlier versions (Lean \texttt{eulerB}) is now only a legacy
special case; nothing below uses product structure.
\end{definition}

The decisive structural property is \textbf{boundedness} of the coefficients. When
\(\|X_n(\omega)\|\le 2\), the tail \(\sum_{n>N}X_n(\omega)n^{-s}\) is dominated by
\(2\sum_{n>N}n^{-a}\) on \(\{\Rea(s)\ge a>1\}\) for \emph{every} \(\omega\), giving
absolute, locally uniform convergence with \emph{no} almost-sure qualifier (a Weierstrass
\(M\)-test; Lean \texttt{bSum\_tail\_small} for the product instance). This is what anchors
the random series to the genuine \(\eta\), whose Dirichlet series converges there.

\begin{definition}[The atomless radius-\(2\) prior]\label{def:prior}
Let \(\texttt{diskLaw}\) be the uniform probability law on the closed disk of radius
\(\mathbf{2}\), \(\overline{B(0,2)}\subset\mathbb{C}\), and
\[
  \texttt{priorBall} \;=\; \bigotimes_{n}\ \texttt{diskLaw}
  \;=\; \texttt{Measure.infinitePi}\,(\lambda\_.\ \texttt{diskLaw})
\]
the independent product over indices. It is a probability measure, supported in the
radius-\(2\) ball (so \(\|X_n\|\le 2\) a.e.), and \textbf{atomless}
(\(\texttt{priorBall}\{\omega\}=0\)). (Lean: \texttt{priorBall\_isProb},
\texttt{priorBall\_ball}, \texttt{priorBall\_atomless}; on disk for radius \(1\), to be
re-proved verbatim at radius \(2\).)
\end{definition}

Two structural properties of the prior carry the argument.

\begin{itemize}
  \item \textbf{Atomlessness} is what prevents any single realization --- in particular
        the deterministic \(\eta\)-realization \(X_n\equiv1\) --- from carrying positive
        mass, so the conditional behaviour at \(X_n\equiv1\) is a genuine
        \emph{neighborhood} (Tjur) limit, not an evaluation.
  \item \textbf{Radius \(2\), not \(1\)} places the mean realization \(X_n\equiv1\) in the
        \emph{open interior} of the support. At radius \(1\), \(X_n\equiv1\) sat on the
        boundary \(\partial\overline{B(0,1)}\) of the support, where the Tjur limit is
        one-sided and degenerate; at radius \(2\) every neighborhood of \(X_n\equiv1\) is a
        genuine two-sided interior neighborhood with positive mass in all directions, so
        the rcp at the conditioning value is unambiguous. Boundedness --- not the value of
        the radius --- is the critical property; radius \(2\) is the smallest convenient
        choice that interiorizes the mean.
\end{itemize}

\subsection{Mean and covariance: choosing the law is choosing the model}\label{sec:general}

The model is specified --- Mehler/Kopperman-style --- by two moments:
\begin{itemize}
  \item \textbf{mean} \(\mathbb{E}[X_n]=1\): the \(X_n\equiv1\) realization is the
        deterministic target series \(\eta\), the conditioning center;
  \item \textbf{covariance} \(\mathbb{E}\bigl[(X-\text{mean})^2\bigr]\): the chosen law,
        which must have \emph{rich enough support} for universality (Section~\ref{sec:nondetect}).
\end{itemize}
The conclusion is a property of \((\text{mean},\text{covariance})\), not of any product or
multiplicative structure (Remark~\ref{rem:nomult}). The radius-\(2\) unit-disk product of
\S\ref{sec:model} is one admissible covariance; a Fock-space picture (first \(N\)
coefficients pinned, the rest free) is another.

\emph{Boundedness is not optional.} The bound \(\|X_n\|\le 2\) gives \emph{absolute},
locally uniform convergence on \(\{\Rea(s)\ge a>1\}\) for \emph{every} \(\omega\), which
anchors the random series to \(\eta\); without it the series need not converge for every
\(\omega\) on \(\Rea>1\) and nothing about \(\eta\) could be concluded. Inside the strip,
\(1/2<\Rea(s)\le1\), absolute convergence fails; there the \(X_n\) being uncorrelated
(independent under the product prior) with bounded variance makes the centered series'
term-variances \(\operatorname{Var}(X_n)\,n^{-2\Rea(s)}\) summable for \(\Rea(s)>1/2\) ---
\textbf{Kolmogorov's one-series theorem} (\(L^2\) from orthogonality, almost-sure via the
\(L^2\)-bounded martingale of partial sums). These two regimes --- absolute on \(\Rea>1\),
Kolmogorov on \(\Rea>1/2\) --- are the right-edge and left-edge regimes of the contour,
and both ingredients are \emph{already in the project}: the deterministic Dirichlet-series
summability, locally-uniform convergence, and analyticity on \(\{\Rea>1/2\}\)
(\texttt{EtaConvergence.lean}: \texttt{summable\_etaPairedTerm},
\texttt{paired\_tendstoUniformlyOn}, \texttt{analyticOnNhd\_paired\_tsum}), and the
i.i.d.\ Mehler product with its independence and variance (\texttt{SphereGaussian.lean}:
\texttt{gammaMeasure}, \texttt{iIndepFun\_infinitePi}, \texttt{variance\_id\_gaussianReal}).
So the convergence/analyticity of the general series on \(\{\Rea>1/2\}\) is an
\emph{assembly} of existing results, not a fresh obligation.

\subsection{The rcp kernel}\label{sec:kernel}

Let \(\partial^\ast R = R.\texttt{closure}\setminus R.\texttt{openInt}\) be the topological
boundary of a rectangle \(R\), and let \(\Lambda\subset\partial^\ast R\) be a \emph{local}
arc that does not enclose \(s_0\) (Remark~\ref{rem:maxmod}). Write the \emph{edge trace}
\(T(\omega) = (z\mapsto \etaX(z,\omega))\restriction \Lambda\) and the \emph{interior
value} \(\etaX(s_0,\omega)\).

\begin{definition}[rcp kernel]\label{def:kernel}
\[
  \texttt{rcpKernel}(s_0,\Lambda)
  \;=\; \texttt{condDistrib}\bigl(\etaX(s_0,\cdot),\ T,\ \texttt{priorBall}\bigr),
\]
the regular conditional distribution of the interior value given the (local) edge trace,
under the bounded prior. (Lean: \texttt{rcpKernel}, built from \texttt{condDistrib};
measurability of both maps and finiteness of \texttt{priorBall} are the inputs.)
\end{definition}

The kernel is well posed because the precondition of Tjur's
characterization~\eqref{eq:tjur} holds: the conditioning value \(e_\eta\) (the trace of the
mean realization) lies in the support of the trace law, with \emph{positive} mass on every
neighborhood --- measured against the deterministic skeleton on which \(\etaX\) centers
(Lean \texttt{priorBall\_edgeNbhd\_pos}, from the positive-mass coefficient box
\texttt{priorBall\_box\_pos}). The radius-\(2\) choice strengthens this from ``in the
support'' to ``in the interior of the support''.

\subsection{The contour side, for free}

On the boundary, reciprocals converge: if \(\etaX(\cdot,\omega_k)\to\eta\) uniformly on
\(\partial^\ast R\) (with \(\eta\) zero-free there, i.e.\ its only strip zero is the
interior \(s_0\)), then
\[
  \oint_{\partial R}\frac{dz}{\etaX(z,\omega_k)}
  \;\longrightarrow\;
  \oint_{\partial R}\frac{dz}{\eta(z)} .
\]
(Proved for the product instance: \texttt{rcp\_recipEulerB\_tendsto\_recipEta}, via a
per-edge dominated-convergence step \texttt{edge\_integral\_tendsto} and continuity
helpers; the strip hypothesis \(\Rea<1\) on the closure makes the integrands continuous.)
For the general series the same locally-uniform convergence is the assembly of
\S\ref{sec:general}. This continuity of evaluation along convergent traces is exactly what
makes the ``zero at \(s_0\)'' functional continuous, and hence the transfer
equality~\eqref{eq:transfer} available; here it is used only as supporting infrastructure
for the probabilistic statements.

\section{The Transfer Equality}\label{sec:transfer}

The single role of the Mehler probability measure is to supply the identity~\eqref{eq:transfer}.

\begin{theorem}[Transfer equality]\label{thm:transfer}
For \(s_0\) with \(1/2<\Rea(s_0)<1\),
\[
  \rcp\bigl(\etaX(s_0)=0 \,\big|\, X_n=1\ \forall n\bigr)
  \;=\; \mathbf{1}\{\eta(s_0)=0\}
  \;=\; \mathbb{P}\bigl(\eta(s_0)=0\bigr).
\]
\end{theorem}
\textbf{(Lean target: \texttt{transfer\_equality}; it replaces the arithmetical bridge
\texttt{counterexample\_iff\_rcpZero} of earlier versions.)}

\paragraph{Why it holds.} The regular conditional probability is the Tjur
limit~\eqref{eq:tjur} at the conditioning value \(e_\eta=\) (the trace of the mean
realization \(X_n\equiv1\)). Three facts make the limit collapse to the deterministic
indicator:
\begin{enumerate}
  \item \emph{The limit exists and is foundation-independent} because the substrate is
        Radon (Polish inner-product space, inner-regular Borel measures) --- the
        Mehler-within-Kopperman layer of Section~\ref{sec:kopperman}.
  \item \emph{The conditioning value is interior} (radius \(2\), \S\ref{sec:model}): every
        neighborhood \(V\downarrow\{T=e_\eta\}\) is a genuine two-sided neighborhood with
        positive mass, so the limit is a single number, not an artifact of a one-sided
        approach.
  \item \emph{The event is continuous at the mean realization.} ``\(\etaX\) has a zero at
        \(s_0\)'' is, by the locally-uniform contour convergence of \S\ref{sec:general},
        continuous in the trace topology at \(X_n\equiv1\); the rcp at an interior support
        point returns the deterministic value of such a functional, namely
        \(\mathbf{1}\{\eta(s_0)=0\}\).
\end{enumerate}
This is the precise content of ``adding the Mehler measure proves that the probability of
a zero of the standard \(\eta\) equals that of \(\etaX\) given \(X_n=1\)''. It is the first
of the route's two deep analytic inputs.

\section{Non-detectability: \texorpdfstring{$L\neq 1$}{L != 1}}\label{sec:nondetect}

\begin{theorem}[No forced zero: \(L\neq 1\)]\label{thm:nondetect}
Conditioned on a \emph{local} edge neighborhood of \(\eta\) that does not enclose \(s_0\),
the conditional law of the interior value \(\etaX(s_0)\) is not a point mass at \(0\):
\[
  L \;=\; \rcp\bigl(\etaX(s_0)=0 \,\big|\, X_n=1\bigr) \;\neq\; 1 .
\]
\end{theorem}
The mechanism is \textbf{Bagchi/Voronin universality}. Among the realizations consistent
with the local edge data there is a positive-mass subset --- an \emph{alternative
average} --- given by a \emph{zero-free} random Dirichlet series that matches \(\eta\) on
the edges but stays nonzero at \(s_0\); universality is precisely the statement that the
prior's support is all zero-free holomorphic functions on a connected-complement compact,
so such a subset exists and keeps \(L<1\). The analytic input is Bagchi (1981); the
closest formal statement available is the universality lemma
\texttt{eulerG\_universality\_threeLeftEdges} of the companion file, to be reused as a
starting point. This is the second of the route's two deep analytic inputs.

\begin{remark}[Local neighborhood, and why not the full boundary]\label{rem:maxmod}
The conditioning must be on a \emph{local} neighborhood of a boundary point --- no
particular shape or size, only that it does not enclose \(s_0\). On the \emph{full}
rectangle boundary the maximum-modulus principle would force any holomorphic realization
\(\varepsilon\)-close to \(\eta\) on \(\partial R\) to be \(\varepsilon\)-close to \(\eta\)
throughout \(R\), hence \(|f(s_0)|<\varepsilon\) --- pinning the interior to
\(\eta(s_0)\) and (if \(\eta(s_0)=0\)) giving \(L=1\). A local, non-enclosing neighborhood
escapes the pin, and is exactly the connected-complement geometry that universality
requires.
\end{remark}

\begin{remark}[Multiplicativity is not needed --- but boundedness is]\label{rem:nomult}
Bagchi-type universality holds for general random Dirichlet series, so the zero-free
alternative average comes from the \emph{full support of the chosen prior law} (its
covariance), not from any Euler-product or multiplicative structure. We use no Euler
product at all. \emph{Boundedness}, by contrast, is \textbf{not} dispensable:
\(\|X_n\|\le 2\) is what gives absolute every-\(\omega\) convergence on \(\Rea>1\) and so
anchors the model to \(\eta\) (\S\ref{sec:general}). The requirements are ``bounded''
(critical) and ``rich enough support'' (universality), not ``multiplicative''.
\end{remark}

\section{Assembly: a Positive RH in ZFC}\label{sec:assembly}

\begin{theorem}[No strip zero]\label{thm:nozero}
For every \(s_0\) with \(1/2<\Rea(s_0)<1\), \(\eta(s_0)\neq 0\).
\end{theorem}
\begin{proof}[Proof (modulo the two inputs)]
Suppose \(\eta(s_0)=0\). By the transfer equality (Theorem~\ref{thm:transfer}),
\(L=\rcp(\etaX(s_0)=0\mid X_n=1)=\mathbf{1}\{\eta(s_0)=0\}=1\). But by non-detectability
(Theorem~\ref{thm:nondetect}), \(L\neq 1\). Contradiction; hence \(\eta(s_0)\neq 0\).
Equivalently, the deterministic indicator \(\mathbf{1}\{\eta(s_0)=0\}\in\{0,1\}\) equals
\(L\neq1\), so it equals \(0\).
\end{proof}

\begin{corollary}[Riemann Hypothesis]\label{cor:rh}
Every nontrivial zero of \(\zeta\) lies on \(\Rea(s)=1/2\).
\end{corollary}
\begin{proof}
By Theorem~\ref{thm:nozero}, \(\eta\) --- hence \(\zeta\), since \(\eta(s)=(1-2^{1-s})
\zeta(s)\) and the factor is zero-free in the open strip --- has no zero with
\(1/2<\Rea(s)<1\); by the functional equation the same holds for \(0<\Rea(s)<1/2\).
\end{proof}

Note what is \emph{absent}: no \(\Pi^0_1\) Schoenfeld sentence, no primitive-recursive
matrix, no \(\Sigma^0_1\) witness extraction, no model/\(T\)-conservativity step, and no
unprovability hedge. The argument lives entirely in standard complex analysis and
elementary measure theory (ZFC); its only non-elementary content is the two analytic
inputs --- the transfer equality and Bagchi universality --- both flagged as the route's
debts.

\section{The Kopperman/Mehler Formalism}\label{sec:kopperman}

This section explains \emph{why} the probabilistic substrate above is the natural home for
the argument, in the language of the \textbf{Kopperman/Mehler formalism} shared with the
companion \(P\neq NP\) development. The punchline: the Mehler measure is the device that
converts the deterministic question ``does \(\eta\) vanish at \(s_0\)?'' into a conditional
probability \(L\) that universality can bound, and the Radon substrate is what makes \(L\)
a well-defined, foundation-independent number.

\subsection{A formalism is a triple}\label{sec:triple}

Following Kopperman (1967) and Mehler (1866), a \emph{formalism} for a piece of analysis is
a triple \(\bigl(\mathcal{H},\ \mathsf{Skel},\ \mu\bigr)\):
\begin{enumerate}
  \item a \textbf{separable Hilbert substrate} \(\mathcal{H}\) --- the ``territory'', whose
        metric completeness (every Cauchy net of approximants has a genuine limit) is the
        analytic ground on which holomorphic limits live;
  \item a \textbf{countable decidable dense skeleton} \(\mathsf{Skel}\subset\mathcal{H}\) ---
        the ``map'', finite combinations with rational data, on which all \emph{computation}
        happens;
  \item an \textbf{atomless probability measure} \(\mu\) on \(\mathcal{H}\) --- the device
        that, since we cannot algorithmically \emph{identify} which skeleton element
        approximates a given target, lets us \emph{average} over them with full measure.
\end{enumerate}
``Choosing the measure \(\mu\) is choosing the model.'' In Lean the triple is the
\texttt{Formalism} structure (\texttt{PnpProof.Kopperman}); the route builds its instance
\texttt{rcpFormalism} from \texttt{formalismOfPrior} applied to the atomless substrate prior.

\subsection{The three layers in this route}\label{sec:layers}

\paragraph{Substrate \(\mathcal{H}\).} The separable Hilbert space
\(L^2(\texttt{priorBall})\) --- \(L^2\) of the infinite product of the uniform radius-\(2\)
disk law. Its existence and that of a basis are \emph{already guaranteed by the Mehler
formalism}, even in this \emph{non-Gaussian} (uniform-on-disk, bounded) case: the Mehler
construction together with the automatic separability of \(L^2\) supply the substrate, with
no Gaussian-specific Hermite/Zernike polynomials and no basis to build by hand. So there is
no ONB-construction, orthonormal-pair, or completeness obligation. What we \emph{do} use is a
\textbf{family of vectors for the rcp limit}: the finite \textbf{indicator tensors}
\(\bigotimes_n \varphi_{\alpha_n}\) (\(\alpha\) finitely supported, \(\varphi_k\) indicator
functions of subsets of the radius-\(2\) disk --- uniform distributions on subsets of the
original ball, trivial factors \(=1\)) that populate the Tjur neighborhoods of \(e_\eta\) and
carry the alternative averages. They need no orthonormality or completeness --- only to be
available as the vectors over which the limit \(L=\lim_V\mathbb{P}(A\cap V)/\mathbb{P}(V)\) is
taken. The one property kept is \textbf{consistency to the original law}: the
finite-coordinate marginals of these indicator tensors reconstruct the original
uniform-on-ball measure (their Kolmogorov extension is exactly \texttt{priorBall}; the
unconstrained coordinates always carry the original \texttt{diskLaw}). The substrate is
\textbf{Radon} --- a Polish inner-product space, Borel measures inner-regular --- which is
what makes the Tjur limit~\eqref{eq:tjur} \emph{well defined and foundation-independent}.
(In Lean: \texttt{SphereGaussian.lean} for the Mehler substrate, \texttt{Measure.infinitePi}
/ \texttt{infinitePi\_pi} for the consistency; the abstract atomless prior already in
\texttt{rcpFormalism} suffices.)

\paragraph{Skeleton \(\mathsf{Skel}\).} The finite Dirichlet approximants and the finite
random partial sums \(\sum_{n\le N}X_n(\omega)n^{-s}\) with rational data: decidable,
computably enumerable, the deterministic skeleton on which \(\etaX\) centers and against
which \texttt{edgeNbhd} positivity is measured.

\paragraph{Measure \(\mu\).} An atomless Mehler-type prior on the coefficients, fixed by its
mean (\(\mathbb{E}[X_n]=1\)) and covariance (\S\ref{sec:general}); the radius-\(2\)
unit-disk product \texttt{priorBall} is one concrete choice. Three properties do the work.
\textbf{Atomlessness} forces the conditioning at \(X_n\equiv1\) to be a neighborhood (Tjur)
phenomenon. \textbf{Interiority of the mean} (radius \(2\)) makes that Tjur limit two-sided
and unambiguous, which is what the transfer equality (Theorem~\ref{thm:transfer}) needs.
\textbf{Full support} (universality of the law) supplies the zero-free \emph{alternative
average} that keeps \(L\neq1\) (Theorem~\ref{thm:nondetect}); this is Bagchi/Voronin and does
not need multiplicativity (Remark~\ref{rem:nomult}). Independently, \textbf{bounded
coefficients} (\(\|X_n\|\le2\)) give the absolute \(\Rea>1\) convergence that anchors the
model to \(\eta\) (\S\ref{sec:general}). ``Choosing the measure is choosing the model'':
different admissible (bounded, full-support) covariances give the same conclusion, because it
is a property of \((\text{mean},\text{covariance})\), not of any one law.

\subsection{What the formalism does, and does not, buy}\label{sec:fences}

\begin{enumerate}
  \item \textbf{Foundation-independence, not a free proof.} The Radon/Mehler structure makes
        the rcp \(L\) a genuine, foundation-independent number (it does not depend on the
        ZFC bookkeeping), and locates the two analytic debts; it does \emph{not} by itself
        prove either of them. Both the transfer equality (Theorem~\ref{thm:transfer}) and the
        limit form of \(L\neq1\) (Theorem~\ref{thm:nondetect}, Bagchi) are genuine analytic
        inputs. Everything else --- the kernel, the prior, the convergence --- is proved or
        assembled from existing results.
  \item \textbf{No new axioms.} The Lean development stays within \texttt{propext},
        \texttt{Classical.choice}, \texttt{Quot.sound}. The Radon substrate and the
        Mehler/bounded prior are the \emph{model in which the probabilistic statement is
        interpreted}, not extra axioms; the final theorem is a standard statement about
        \(\zeta\).
  \item \textbf{A positive theorem, not an unprovability one.} Unlike the earlier
        \(\Pi^0_1\) encoding --- which could conclude only that a counterexample is
        unwitnessable --- the ZFC-direct argument turns \(L\neq1\) into the positive
        \(\eta(s_0)\neq0\) via the two-valuedness of the deterministic indicator. There is
        no remaining gap between ``unprovable'' and ``true''.
\end{enumerate}

\noindent In short: the route is the Kopperman/Mehler formalism with \(\mathcal{H}=\) the
Radon substrate, \(\mathsf{Skel}=\) finite Dirichlet approximants, and \(\mu=\) the atomless
radius-\(2\) prior fixed by its mean and covariance; the Mehler measure converts ``does
\(\eta\) vanish at \(s_0\)?'' into the conditional probability \(L\), the transfer equality
identifies \(L\) with the deterministic indicator, and universality keeps \(L\neq1\) --- so
the indicator is \(0\) and \(\eta(s_0)\neq0\). Two deep inputs (transfer equality, Bagchi),
everything else \texttt{sorry}-free or assembled probability.

\section{Lean 4 Formalization Status}\label{sec:status}

The probabilistic core is formalized in \texttt{RcpEuler.lean}, imported by the build root
\texttt{RiemannProof.lean}. The plan is \texttt{IMPLEMENTATION\_PLAN\_RCP.md}.

\paragraph{On disk (sorry-free), for the earlier radius-\(1\) Euler-product instance.}
\(\mathcal{B}_P\) never-zero and analytic (\texttt{eulerB\_ne\_zero},
\texttt{eulerB\_differentiableOn}); the bounded-coefficient \(M\)-test
(\texttt{bSum\_tail\_small}); the atomless prior (\texttt{diskLaw}, \texttt{priorBall},
\texttt{priorBall\_isProb}, \texttt{priorBall\_ball}, \texttt{priorBall\_atomless}); the
\texttt{condDistrib} kernel (\texttt{rcpKernel}); the positive-mass Tjur precondition
(\texttt{priorBall\_edgeNbhd\_pos}); and the reciprocal contour convergence
(\texttt{rcp\_recipEulerB\_tendsto\_recipEta}). The general
\((\text{mean},\text{covariance})\) convergence/analyticity on \(\{\Rea>1/2\}\) reuses
existing project results (\texttt{EtaConvergence.lean}, \texttt{SphereGaussian.lean}), so it
is an assembly. There is no substrate-construction obligation: the Mehler formalism already
guarantees the separable substrate and a basis (non-Gaussian case included), and the
abstract atomless prior in \texttt{rcpFormalism} suffices.

\paragraph{The redesign of this paper (not yet on disk).} Three changes retarget the core,
detailed in \texttt{IMPLEMENTATION\_PLAN\_RCP.md} (master section ``The 2026-06-22
redesign''):
\begin{enumerate}
  \item[\textbf{(R)}] move \texttt{diskLaw} to the radius-\(2\) disk (\texttt{closedBall 0
        2}); re-prove \texttt{priorBall\_*} (mechanical), interiorizing the mean realization;
  \item[\textbf{(M)}] generalize \(\mathcal{B}_P\) to the general, non-multiplicative random
        Dirichlet series \(\etaX=\sum_n X_n n^{-s}\) (\(\|X_n\|\le2\), mean \(1\));
  \item[\textbf{(Z)}] replace the arithmetical \texttt{SchoenfeldPRA.lean} spine --- the
        primitive-recursive Schoenfeld matrix, the \(\Pi^0_1\) interpreter, the bridge
        \texttt{counterexample\_iff\_rcpZero}, and \(\texttt{RH\_PRA}\) --- by the
        measure-theoretic \texttt{transfer\_equality} and a final standard-RH assembly. The
        Schoenfeld/\(\Pi^0_1\) material is \emph{dropped}, not discharged.
\end{enumerate}

\paragraph{The two genuine debts.}
\begin{itemize}
  \item \textbf{The transfer equality} \(\rcp(\etaX(s_0)=0\mid X_n=1)=\mathbf{1}\{\eta(s_0)
        =0\}\) (Theorem~\ref{thm:transfer}; Lean \texttt{transfer\_equality}). Its content is
        that the rcp at the interior mean realization recovers the deterministic value of the
        continuous ``zero at \(s_0\)'' functional. It is the new bridge, replacing the
        arithmetical one.
  \item \textbf{Non-detection} \(L\neq1\) (Theorem~\ref{thm:nondetect}; Lean
        \texttt{not\_rcpZeroAt}, limit form on a local neighborhood). The on-disk
        \texttt{not\_rcpZeroAt} proves only a trivial \(\forall k\) placeholder; the
        substantive \(L\neq1\) is \textbf{Bagchi (1981) universality}, with the closest
        existing formal statement \texttt{eulerG\_universality\_threeLeftEdges}.
\end{itemize}
Both are flagged as \texttt{sorry} with recipes; neither is to be improvised.

\section*{Status summary}

The probabilistic skeleton --- the random Dirichlet series, the \emph{atomless} prior, the
regular conditional kernel, the positive-mass Tjur precondition, and the reciprocal contour
convergence --- is formalized and \texttt{sorry}-free for a radius-\(1\) Euler-product
instance (\texttt{RcpEuler.lean}), with the general \((\text{mean},\text{covariance})\)
convergence reusing existing project results (\S\ref{sec:general}). The redesign to
radius~\(2\), the general non-multiplicative series, and the measure-theoretic transfer
equality in place of the arithmetical bridge is the current plan and not yet on disk. The
genuine debts are two deep, un-formalized analytic inputs: the \textbf{transfer equality}
(\(\rcp(\etaX(s_0)=0\mid X_n=1)=\mathbf{1}\{\eta(s_0)=0\}\), the role of the Mehler measure)
and \textbf{Bagchi (1981) universality} (the local-neighborhood non-detection \(L\neq1\)).
Read through the Kopperman/Mehler formalism, the route says: the Mehler measure turns ``does
\(\eta\) vanish at \(s_0\)?'' into a conditional probability \(L\); the transfer equality
identifies \(L\) with the deterministic \(\{0,1\}\) indicator; and conditioning on a
\emph{local} edge neighborhood, universality keeps \(L\neq1\) --- so the indicator is \(0\)
and \(\eta(s_0)\neq0\), uniformly over the strip, which is RH. This is a positive analytic
conclusion in ZFC, not an unprovability statement. Until the two inputs are discharged, the
development is a \emph{verified reduction} of the Riemann Hypothesis to them, not a complete
proof.

\end{document}
